\documentclass[11pt]{article}
\usepackage[margin=1in]{geometry}
\usepackage{enumitem}
\usepackage{xcolor}
\usepackage{amsmath,amssymb}
\usepackage[backend=biber,style=numeric,url=false,isbn=false,doi=false]{biblatex}
\addbibresource{biblio_ConfESC.bib}
\usepackage{titlesec}

% Adjust heading sizes here:
\titleformat{\part}[display]{\centering\bfseries\large}{}{0pt}{}
\titleformat{\section}{\bfseries\large}{}{0pt}{}
\titleformat{\subsection}{\bfseries\normalsize}{}{0pt}{}

\newcommand{\referee}[1]{\textbf{Referee:} \textit{#1}}
\newcommand{\response}[1]{\textbf{Response:} #1}
\newcommand{\placeholder}[1]{{\color{red}\textbf{[TODO: #1]}}}

\begin{document}

\begin{center}
  {\Large\bfseries Response to Referee Reports}\\[6pt]
  {\large On exponential separation of analytic self-conformal sets on the real line}\\[4pt]
  B\'ar\'any, Kolosv\'ary, and Troscheit
\end{center}

\bigskip

\noindent We thank both referees for their careful reading of our paper and their helpful suggestions.
Below we address all comments point by point.  All requested corrections have been implemented in
the revised manuscript.

\bigskip

\part*{Response to Referee X}

\section*{Main issue}

\referee{The only detail I did not fully understand occurs in the proof of Theorem 1.12, where the
authors apply Lemma 4.1, which requires the IFS to have no exact overlaps, but you pick an arbitrary
IFS. I may be misunderstanding something here, but it would be helpful if the authors clarified this
point.}

\response{Thank you for bringing this to our attention. We believe the referee is referring to Theorem~1.4,
where Lemma~4.1 is applied in the density part of the proof. (This appears on the page referenced.
Theorem 1.12 itself does not use Lemma 4.1)

Since we are proving denseness, we may assume without loss of
generality that the given IFS~$\Phi$ has no exact overlaps: if it does, an arbitrarily small
perturbation removes them, and the argument then proceeds for this perturbed system. We have added
an explicit sentence to the proof to clarify this point.}

\section*{Minor corrections}

\begin{itemize}[leftmargin=*]

\item \textbf{1.} Added explicit definition of $\mathcal{B}_\delta$ in condition~(A).

\item \textbf{2.} Removed trailing comma in the $m>n$ case of $\mathbf{i}_m^n$.

\item \textbf{3.} Changed semicolon to colon after Corollary~1.7.

\item \textbf{4, 5, 6, 9, 10, 15, 17, 19, 20, 22, 23, 24, 27, 28, 30.}
All internal cross-references use \texttt{\textbackslash cref\{\}}, and there was a bug in the
compilation when these labels were generated. These issues have now been resolved.

\item \textbf{7.} Redefined $c_{\min}$ and $c_{\max}$ globally over
$\overline{\mathcal{B}_\epsilon}$ instead of $[0,1]$, so the bounds hold in the complex
neighbourhood. The example in Section~1.2.1 needed a small update with this change.

\item \textbf{8.} Added missing commas around cross-references.

\item \textbf{11.} Fixed denominator to $(f_i'(x))^2$.

\item \textbf{12, 14.} Added the missing factor $|\mathbf{i}|$ in Eq.~(2.5) and before Eq.~(2.6).

\item \textbf{13.} Changed $c_{\min}$ to $c_{\min}^2$ before Eq.~(2.6).

\item \textbf{16.} Changed the final equality to $\leq$ in the proof of Lemma~2.8.

\item \textbf{18.} Added explicit statement that $\eta_{n_\ell}<1/4$ for all $\ell$.

\item \textbf{21.} Renamed the enumeration index from $k$ to $r$ in the proof of Lemma~4.1.

\item \textbf{26.} Changed $h$ to $k$ in $(F_{\mathbf{i}}h)(x_{\mathbf{i},\mathbf{j}})$.

\item \textbf{29.} Clarified the passage to read: ``The assertion~(b) follows by applying~(a) to
the sub-IFS $(f_{\mathbf{i}}, f_{\mathbf{j}})$, so we finish the proof by showing~(a).''

\end{itemize}


\part*{Response to Referee Y}

\referee{I suggest the authors take a look at the paper arXiv:2501.13729 by Shunsuke Usuki, which may be relevant.}

\response{We thank the referee for this reference. Usuki~\cite{Usuki_LqMobius25arXiv} proves a
dichotomy for the $L^q$ spectrum of stationary measures for M\"obius IFSs satisfying the strongly
Diophantine condition: either the $L^q$ dimension takes the expected value for all $q>1$, or the
spectrum eventually becomes linear. A key observation is that the latter case occurs whenever two
maps in the IFS share a common fixed point in the attractor. This is directly relevant to our work,
as it provides the general mechanism behind the $L^q$ dimension drop in the example of
Section~1.2.1. We have added a discussion of Usuki's result in Section~1.2 (after Corollary~1.7),
where we already note that our typicality result does not extend to $L^q$ dimensions.}




\section*{Additional changes since submission}

In addition to the corrections suggested by the referee, the following changes have been made to the
manuscript since the version submitted in September 2025:

\begin{itemize}
  \item We added an acknowledgement thanking Amir Algom for pointing out useful references on
    conjugation of analytic IFSs to self-similar IFSs.

  \item In Section~1.2.3 (Conjugation to self-similar systems), we added a reference to Algom,
    Rodriguez Hertz, and Wang \cite{AlgomHertzWang23} and their dichotomy result (Claim~6.1 of that
    paper), noting that our characterisation is complementary.

  \item We added a paragraph on the uniform non-integrability (UNI) condition of Dolgopyat and
    Chernov, explaining its relation to our dual IFS framework: the UNI condition is equivalent to
    $\inf_{x\in I} |H_{\mathbf{i}}(x) - H_{\mathbf{j}}(x)| > 0$ for some
    $\mathbf{i},\mathbf{j}\in\Sigma$, which is strictly stronger than the negation of the
    conjugation characterisation in Theorem~1.14.

  \item Several minor typos (``the there'' $\to$ ``there'', ``Similar question'' $\to$ ``Similar
    questions'', ``and by Using'' $\to$ ``and by using'', trailing comma removal) were fixed
    independently.
\end{itemize}


\printbibliography

\end{document}
